% =====================================================================
% Clase -- Álgebra 9° -- Trimestre I -- Semana 04
% Sesión 007 (mar 22 sep., 13:00, 40 min) y sesión 008 (vie 25 sep., 6:55, 55 min)
% Tema: Números reales
% Subtemas: Radicales: raíz enésima y simplificación / Suma y resta de radicales semejantes
% Hilo conductor del trimestre I: «Al-Juarismi y el nacimiento del álgebra»
% Tema 1 de la guía: Números reales (tema:numerosreales)
% Material del docente (no se entrega a los estudiantes).
% =====================================================================
\documentclass[11pt]{article}
\makeatletter\def\input@path{{../../../../../plantillas/estilo/}}\makeatother
\usepackage{mmcantillo}
\guiadelestudiante{../../guia-didactica/guia-periodo-I-reales-y-factorizacion}

\tipodocumento{Clase}
\asignatura{Álgebra}
\titulo{Semana 4: radicales que se pueden juntar}
\grado{9°}
\periodo{I}

% DBA: matematicas grado 9 · DBA 1 — Utiliza los números reales (sus operaciones, relaciones
% y propiedades) para resolver problemas con expresiones polinómicas. Evidencias de estas dos
% sesiones: «Construye representaciones geométricas y numéricas de los números reales (con
% decimales, raíces, razones, y otros símbolos) y realiza conversiones entre ellas» y, en la
% sesión 008, «Identifica la diferencia entre exactitud y aproximación».

\begin{document}

\encabezado*

\begin{center}
{\renewcommand{\arraystretch}{1.3}
\begin{tabularx}{\linewidth}{@{}l l l X l@{}}
  \toprule
  \textbf{Sesión} & \textbf{Fecha} & \textbf{Min} & \textbf{Subtema} & \textbf{Entregables}\\
  \midrule
  007 & mar 22 sep., 13:00 & 40 & Radicales: raíz enésima y simplificación & ---\\
  008 & vie 25 sep., 6:55 & 55 & Suma y resta de radicales semejantes & Taller y tarea\\
  \bottomrule
\end{tabularx}}
\end{center}

Segunda semana del tema 1. Al-Juarismi operaba con raíces sin tener nuestros símbolos: las
escribía con palabras. Estas dos sesiones son justamente la herramienta que a él le faltaba —
poder escribir $5\sqrt{2}$ y saber que dos de esos se suman como se suman dos manzanas.
Referencia: \refguia{tema:numerosreales}.

% ---------------------------------------------------------------------
\seccion{Sesión 007 — martes 22 de septiembre · 13:00 · 40 min}

\textbf{Objetivo.} Leer y escribir raíces enésimas, y simplificar un radical extrayendo los
factores que son potencias perfectas.

\textbf{Materiales.} Tablero; la guía. No hace falta calculadora (y conviene que no la usen).

\begin{center}
{\renewcommand{\arraystretch}{1.3}
\begin{tabularx}{\linewidth}{@{}l l X@{}}
  \toprule
  \textbf{Momento} & \textbf{Min} & \textbf{Actividad}\\
  \midrule
  Revisión & 0--5 & Revisar la tarea de la semana pasada. Tomar una respuesta con
    $\sqrt{18}$ sin simplificar como puente al tema de hoy.\\
  Notación & 5--14 & $\sqrt[n]{a}$: índice, radicando, raíz. Leerlas en voz alta
    ($\sqrt[3]{8}$, «raíz cúbica de 8»). Por qué $\sqrt{a^2} = |a|$ y no siempre $a$: probar
    con $a = -3$. Es el detalle que más adelante arruina ecuaciones.\\
  Simplificar & 14--32 & La idea: buscar dentro del radical un factor que sea potencia
    perfecta y sacarlo. $\sqrt{18} = \sqrt{9 \cdot 2} = 3\sqrt{2}$;
    $\sqrt{75} = 5\sqrt{3}$; $\sqrt[3]{54} = 3\sqrt[3]{2}$. Método al tablero:
    descomponer en factores primos y agrupar de a $n$. Hacer dos con el curso y dos solos.\\
  Cierre & 32--40 & Pregunta que deja abierta la sesión 008: «¿cuánto es
    $\sqrt{18} + \sqrt{2}$?» Dejar que tanteen; casi seguro alguien dirá $\sqrt{20}$.
    No corregir todavía: mañana se ve por qué no.\\
  \bottomrule
\end{tabularx}}
\end{center}

% ---------------------------------------------------------------------
\seccion{Sesión 008 — viernes 25 de septiembre · 6:55 · 55 min}

\textbf{Objetivo.} Sumar y restar radicales semejantes, reconocer cuándo \emph{no} se pueden
juntar, y usarlo para calcular perímetros exactos.

\textbf{Materiales.} Tablero; regla; la guía.

\begin{center}
{\renewcommand{\arraystretch}{1.3}
\begin{tabularx}{\linewidth}{@{}l l X@{}}
  \toprule
  \textbf{Momento} & \textbf{Min} & \textbf{Actividad}\\
  \midrule
  La trampa & 0--10 & Retomar $\sqrt{18} + \sqrt{2}$. Comprobar numéricamente:
    $\num{4,243} + \num{1,414} = \num{5,657}$, mientras que $\sqrt{20} = \num{4,472}$. No
    son iguales. Entonces, ¿cuánto es? Simplificando: $3\sqrt{2} + \sqrt{2} = 4\sqrt{2}$.
    \textbf{La raíz no se reparte en la suma.}\\
  Semejantes & 10--24 & Radicales semejantes: mismo índice y mismo radicando \emph{después
    de simplificar}. Se suman los coeficientes, como $3x + x = 4x$. Ejemplos de la guía:
    $5\sqrt{27} - \sqrt{147} + \sqrt{12} = 10\sqrt{3}$. Y uno que no se puede:
    $\sqrt{2} + \sqrt{3}$ se queda así.\\
  Perímetros & 24--34 & Aplicación: el perímetro de un triángulo rectángulo isósceles de
    catetos 5 es $10 + 5\sqrt{2}$. Exacto contra aproximado: $\num{17,07}$ es una
    aproximación; $10 + 5\sqrt{2}$ es el número. Volver a la distinción del tema.\\
  Taller & 34--50 & Taller individual (5 puntos).\\
  Cierre & 50--55 & Recoger el taller y dejar la tarea. Anuncio: la próxima semana,
    multiplicar y dividir radicales — ahí sí la raíz se reparte, y conviene ver por qué.\\
  \bottomrule
\end{tabularx}}
\end{center}

\begin{nota}
\textbf{El error que hay que cazar.} $\sqrt{a + b} \neq \sqrt{a} + \sqrt{b}$. Vale la pena
dejarlo escrito en el tablero toda la semana, con el contraejemplo al lado:
$\sqrt{9 + 16} = 5$ pero $\sqrt{9} + \sqrt{16} = 7$. Es el mismo error que reaparecerá en
potencias y en logaritmos.
\end{nota}

\begin{nota}
\textbf{Pendiente.} Los ejercicios de esta semana vienen del banco de Álgebra 8°
(\texttt{numeros-reales-8}, \texttt{medidas-con-radicales-8}): falta decidir cómo se reparten
los irracionales entre 8° y 9° para no repetir el trimestre entero (ver \texttt{PENDIENTES.md}).
\end{nota}

\end{document}
